\worksheettitle
  {M07-R. От условия к полному списку}
  {\modulemark{M07-R} Коррекционная карточка}

\studentnameblock

\begin{notebox}[Что нужно получить]
  Нужно не просто найти несколько чисел, а получить \textbf{все} подходящие
  числа. Для этого выберите одну цифру, проверьте по порядку все её возможные
  значения и каждый раз вычисляйте вторую цифру.
\end{notebox}

\begin{practicebox}[Шаг 1. Запишите условие]
  Цифра единиц на \(4\) больше цифры десятков.

  Цифра десятков --- \(a\), цифра единиц --- \(b\).

  Больше цифра \answerblank[22mm]. Поэтому равенство:
  \answerblank[35mm].
\end{practicebox}

\begin{practicebox}[Шаг 2. Выберите цифру для перебора]
  Из равенства \(b=a+4\) удобно сначала выбирать \(a\), а затем вычислять
  кандидат \(a+4\) на значение \(b\).

  \begin{center}
  \renewcommand{\arraystretch}{1.5}
  \begin{tabular}{c|ccccccccc}
    \toprule
    перебираем \(a\) & 1 & 2 & 3 & 4 & 5 & 6 & 7 & 8 & 9 \\
    \midrule
    кандидат \(a+4\) & 5 & \answerblank[7mm] & \answerblank[7mm] &
      \answerblank[7mm] & \answerblank[7mm] & 10 & 11 & 12 & 13 \\
    \bottomrule
  \end{tabular}
  \end{center}

  Обведите кандидаты от \(0\) до \(9\). Остальные значения зачеркните: они не
  могут стать цифрой \(b\).
\end{practicebox}

\begin{practicebox}[Шаг 3. Составьте числа]
  Пара \((a,b)=(1,5)\) даёт число \(15\).

  Полный список: \answerblank[95mm].

  Какие значения \(a\) вы проверили? \answerblank[65mm].

  Почему список полный? \writinglines{1}
\end{practicebox}

\newpage
\begin{notebox}[Иногда удобнее перебирать \(b\)]
  Если условие записано как \(a=2b\), по выбранному \(b\) сразу вычисляется
  цифра \(a\). Поэтому здесь удобнее перебирать \(b\), а не \(a\).
\end{notebox}

\begin{practicebox}[Шаг 4. Цифра десятков равна удвоенной цифре единиц]
  Равенство: \(a=2b\).

  \begin{center}
  \renewcommand{\arraystretch}{1.5}
  \begin{tabular}{c|cccccc}
    \toprule
    перебираем \(b\) & 0 & 1 & 2 & 3 & 4 & 5 \\
    \midrule
    кандидат \(2b\) & 0 & \answerblank[9mm] & \answerblank[9mm] &
      \answerblank[9mm] & \answerblank[9mm] & 10 \\
    \bottomrule
  \end{tabular}
  \end{center}

  При \(b=0\) получается \(a=0\), но цифра десятков двузначного числа не может
  быть нулём. При \(b=5\) получается \(a=10\), а \(10\) не цифра. При
  \(b>5\) значение \(2b\) ещё больше, поэтому допустимой цифры \(a\) также
  не получится.

  Полный список: \answerblank[85mm].

  Почему других чисел нет? \writinglines{1}
\end{practicebox}

\begin{checkpointbox}[Повторная проверка]
  Цифра единиц на \(2\) меньше цифры десятков.
  \begin{enumerate}[label=\textbf{\arabic*.}]
    \item Математическая запись: \answerblank[38mm].
    \item Перебираю цифру: \answerblank[18mm].
    \item Полный список: \answerblank[90mm].
    \item Какие значения проверены и почему список полный?
      \writinglines{2}
  \end{enumerate}
\end{checkpointbox}
